\documentclass[12pt, a4paper]{article}
\usepackage[UTF8]{ctex}
\usepackage{geometry}
\usepackage{listings}
\usepackage{xcolor}
\usepackage{hyperref}
\usepackage{enumitem}
\usepackage{titlesec}
\usepackage{booktabs}
\usepackage{array}
\usepackage{fancyhdr}
\usepackage{graphicx}
\usepackage{tcolorbox}
\tcbuselibrary{listingsutf8, skins}

\geometry{top=2.5cm, bottom=2.5cm, left=3cm, right=3cm}

% 代码样式定义
\definecolor{codebg}{RGB}{245,245,245}
\definecolor{codeframe}{RGB}{200,200,200}
\definecolor{keyword}{RGB}{0,0,200}
\definecolor{comment}{RGB}{0,128,0}
\definecolor{string}{RGB}{163,21,21}
\definecolor{awsorange}{RGB}{255,153,0}
\definecolor{awsdark}{RGB}{35,47,62}
\definecolor{noteblue}{RGB}{232,244,253}
\definecolor{noteframe}{RGB}{33,150,243}
\definecolor{warnbg}{RGB}{255,243,224}
\definecolor{warnframe}{RGB}{255,152,0}

\lstset{
  backgroundcolor=\color{codebg},
  basicstyle=\ttfamily\small,
  keywordstyle=\color{keyword}\bfseries,
  commentstyle=\color{comment},
  stringstyle=\color{string},
  frame=single,
  rulecolor=\color{codeframe},
  breaklines=true,
  showstringspaces=false,
  numbers=left,
  numberstyle=\tiny\color{gray},
  numbersep=5pt,
  tabsize=2,
}

% 页眉页脚
\pagestyle{fancy}
\fancyhf{}
\fancyhead[L]{\textcolor{awsdark}{\textbf{AWS Jam Challenge}}}
\fancyhead[R]{\textcolor{awsorange}{\textbf{Task 1 总结}}}
\fancyfoot[C]{\thepage}
\renewcommand{\headrulewidth}{1pt}

% 标题格式
\titleformat{\section}{\large\bfseries\color{awsdark}}{}{0em}{}[\titlerule]
\titleformat{\subsection}{\normalsize\bfseries\color{awsdark}}{}{0em}{}

\begin{document}

% 标题页
\begin{center}
  {\Huge \textbf{\textcolor{awsdark}{AWS Jam Challenge}}}\\[0.5em]
  {\Large \textcolor{awsorange}{\rule{10cm}{2pt}}}\\[0.5em]
  {\LARGE \textbf{Task 1：启用 RDS MySQL 审计日志并导出至 CloudWatch}}\\[1em]
  {\normalsize \textcolor{gray}{知识点总结 \quad|\quad 操作步骤 \quad|\quad 核心概念}}\\[0.5em]
  {\normalsize \textcolor{gray}{\today}}
\end{center}

\vspace{1em}
\noindent\rule{\textwidth}{0.4pt}
\vspace{1em}

%-----------------------------------------------
\section{任务背景}
%-----------------------------------------------

\begin{tcolorbox}[
  colback=noteblue,
  colframe=noteframe,
  title={\textbf{任务描述}},
  fonttitle=\bfseries
]
Amazon RDS MySQL 实例 \texttt{challenge-node-777} 未启用审计日志（Audit Logging / Option Group），
导致无法对数据库操作进行追踪与取证。\\

\textbf{目标：}在 RDS MySQL 上启用审计日志（通过 Option Group 配置
\texttt{MARIADB\_AUDIT\_PLUGIN}），并将审计日志推送至 Amazon CloudWatch。
\end{tcolorbox}

\vspace{0.5em}

\subsection{涉及服务}

\begin{center}
\begin{tabular}{>{\bfseries}lll}
\toprule
服务名称 & 作用 & 备注 \\
\midrule
Amazon RDS MySQL & 数据库实例 & \texttt{challenge-node-777} \\
Option Group & 插件配置容器 & 承载 \texttt{MARIADB\_AUDIT\_PLUGIN} \\
MARIADB\_AUDIT\_PLUGIN & 审计插件 & 记录连接与查询事件 \\
Amazon CloudWatch & 日志存储与查询平台 & 接收审计日志流 \\
\bottomrule
\end{tabular}
\end{center}

%-----------------------------------------------
\section{核心知识点}
%-----------------------------------------------

\subsection{什么是 Option Group？}

Option Group 是 Amazon RDS 的一种配置容器，用于为数据库引擎启用\textbf{额外功能插件}。
每个 RDS 实例可关联一个 Option Group，通过在其中添加 Option 来扩展数据库能力。

\begin{tcolorbox}[colback=codebg, colframe=codeframe, title=Option Group 关键概念]
\begin{itemize}[leftmargin=1em]
  \item 每个 Option Group 绑定特定的\textbf{数据库引擎}和\textbf{版本}（如 MySQL 8.0）
  \item 一个实例同时只能关联\textbf{一个} Option Group
  \item 修改 Option Group 后，部分选项需要\textbf{重启实例}才能生效
  \item AWS 提供的 \texttt{default} Option Group \textbf{不可修改}，需创建自定义 Option Group
\end{itemize}
\end{tcolorbox}

\subsection{什么是 MARIADB\_AUDIT\_PLUGIN？}

\texttt{MARIADB\_AUDIT\_PLUGIN} 是 RDS MySQL 支持的审计插件，
用于记录数据库的\textbf{连接事件}与\textbf{查询事件}，并将日志写入文件或 CloudWatch。

\begin{center}
\begin{tabular}{lll}
\toprule
参数名 & 说明 & 推荐值 \\
\midrule
\texttt{SERVER\_AUDIT\_LOGGING} & 是否开启审计 & \texttt{ON} \\
\texttt{SERVER\_AUDIT\_EVENTS} & 记录的事件类型 & \texttt{CONNECT, QUERY} \\
\texttt{SERVER\_AUDIT\_QUERY\_LOG\_LIMIT} & 单条查询日志长度限制（字节）& \texttt{1024} \\
\texttt{SERVER\_AUDIT\_INCL\_USERS} & 仅审计指定用户（空=全部）& 按需设置 \\
\texttt{SERVER\_AUDIT\_EXCL\_USERS} & 排除指定用户 & 可排除 \texttt{rdsadmin} \\
\bottomrule
\end{tabular}
\end{center}

\vspace{0.5em}
\begin{tcolorbox}[colback=warnbg, colframe=warnframe, title={\textbf{注意}}]
\texttt{SERVER\_AUDIT\_EVENTS} 需要包含能记录 SQL 文本的事件类型，例如 \texttt{QUERY} 或 \texttt{QUERY\_DML}，才能捕获 DML 操作。\\
若仅设置 \texttt{CONNECT}，则只记录登录/登出事件，无法记录 SQL 语句。
\end{tcolorbox}

\subsection{CloudWatch 日志导出}

RDS 支持将以下类型的日志直接推送至 CloudWatch Logs：

\begin{itemize}
  \item \textbf{Audit log}（审计日志）$\leftarrow$ 本任务目标
  \item Error log（错误日志）
  \item General log（通用日志）
  \item Slow query log（慢查询日志）
\end{itemize}

日志组命名规则：
\begin{lstlisting}[language=bash]
/aws/rds/instance/<实例名>/audit
# 本任务对应：
/aws/rds/instance/challenge-node-777/audit
\end{lstlisting}

%-----------------------------------------------
\section{操作步骤}
%-----------------------------------------------

\subsection{Step 1：创建自定义 Option Group}

\begin{enumerate}[leftmargin=2em]
  \item 登录 \textbf{AWS 管理控制台}，进入 \textbf{RDS} 服务
  \item 在左侧菜单点击 \textbf{Option groups}
  \item 点击 \textbf{Create group}，填写以下信息：
  \begin{itemize}
    \item \textbf{Name}：自定义名称，如 \texttt{my-audit-option-group}
    \item \textbf{Description}：描述信息
    \item \textbf{Engine}：\texttt{MySQL Community}
    \item \textbf{Major engine version}：与实例版本一致（如 \texttt{8.0}）
  \end{itemize}
  \item 点击 \textbf{Create} 完成创建
\end{enumerate}

\subsection{Step 2：添加 MARIADB\_AUDIT\_PLUGIN 选项}

\begin{enumerate}[leftmargin=2em]
  \item 进入刚创建的 Option Group，点击 \textbf{Add option}
  \item 在 \textbf{Option name} 中选择 \texttt{MARIADB\_AUDIT\_PLUGIN}
  \item 配置关键参数：
  \begin{itemize}
    \item \texttt{SERVER\_AUDIT\_LOGGING} = \texttt{ON}
    \item \texttt{SERVER\_AUDIT\_EVENTS} = \texttt{CONNECT, QUERY}
    \item \texttt{SERVER\_AUDIT\_QUERY\_LOG\_LIMIT} = \texttt{1024}
  \end{itemize}
  \item 点击 \textbf{Add option} 保存
\end{enumerate}

\subsection{Step 3：将 Option Group 关联到 RDS 实例}

\begin{enumerate}[leftmargin=2em]
  \item 进入 \textbf{RDS → Databases}，点击实例 \texttt{challenge-node-777}
  \item 点击右上角 \textbf{Modify}（修改）
  \item 在 \textbf{Additional configuration} 部分找到 \textbf{Option group}
  \item 从下拉菜单中选择刚创建的 Option Group
  \item 继续向下找到 \textbf{Log exports} 部分，勾选 \textbf{Audit log}
  \item 点击 \textbf{Continue}，选择 \textbf{Apply immediately}（立即应用）
  \item 点击 \textbf{Modify DB instance} 完成修改
\end{enumerate}

\begin{tcolorbox}[colback=warnbg, colframe=warnframe, title={\textbf{补充说明}}]
为 MySQL 添加 \texttt{MARIADB\_AUDIT\_PLUGIN} 后，审计通常会在 Option Group 生效后立即开始，不需要手动重启实例；但 AWS 官方仍提示该变更\textbf{可能造成短暂中断}，生产环境建议安排在维护窗口或低峰时段执行。
\end{tcolorbox}

\subsection{Step 4：验证日志是否推送至 CloudWatch}

\begin{enumerate}[leftmargin=2em]
  \item 进入 \textbf{CloudWatch → Logs → Log groups}
  \item 搜索日志组：
\begin{lstlisting}[language=bash]
/aws/rds/instance/challenge-node-777/audit
\end{lstlisting}
  \item 若日志组存在且有数据流入，则说明配置成功
\end{enumerate}

%-----------------------------------------------
\section{审计日志格式说明}
%-----------------------------------------------

审计日志每行为 CSV 格式，字段含义如下：

\begin{lstlisting}
20260309 10:37:00,ip-172-17-5-73,rdsadmin,localhost,8,493,QUERY,mysql,'SELECT 1',0,,
\end{lstlisting}

\begin{center}
\begin{tabular}{cll}
\toprule
字段位置 & 字段名 & 示例值 \\
\midrule
1 & 时间戳 & \texttt{20260309 10:37:00} \\
2 & 服务器主机名 & \texttt{ip-172-17-5-73} \\
3 & \textbf{用户名（USERNAME）} & \texttt{rdsadmin} \\
4 & 客户端主机 & \texttt{localhost} \\
5 & 连接 ID & \texttt{8} \\
6 & 查询 ID & \texttt{493} \\
7 & 操作类型 & \texttt{QUERY} / \texttt{CONNECT} \\
8 & 数据库名 & \texttt{mysql} \\
9 & SQL 语句 & \texttt{'SELECT 1'} \\
10 & 返回码 & \texttt{0}（成功） \\
\bottomrule
\end{tabular}
\end{center}

%-----------------------------------------------
\section{常见问题排查}
%-----------------------------------------------

\begin{tcolorbox}[colback=warnbg, colframe=warnframe, title={\textbf{问题1：CloudWatch 日志组中只有 SELECT，没有 UPDATE 记录}}]
\textbf{原因：}审计日志只记录启用之后的操作，历史操作不会被追溯。\\
\textbf{解决：}等待攻击行为再次触发，或手动执行一条 UPDATE 语句进行测试。
\end{tcolorbox}

\vspace{0.5em}

\begin{tcolorbox}[colback=warnbg, colframe=warnframe, title={\textbf{问题2：CloudWatch Insights 查询返回 0 条结果}}]
\textbf{可能原因：}
\begin{itemize}
  \item 时间范围选择过短（默认 1 小时），应调整为更长时间范围
  \item 选错了日志组，注意区分 \texttt{/audit}、\texttt{/error}、\texttt{/general}
  \item \texttt{SERVER\_AUDIT\_EVENTS} 未包含 \texttt{QUERY}，导致 DML 未被记录
\end{itemize}
\end{tcolorbox}

%-----------------------------------------------
\section{总结}
%-----------------------------------------------

\begin{tcolorbox}[colback=noteblue, colframe=noteframe, title={\textbf{Task 1 核心要点}}]
\begin{enumerate}
  \item RDS MySQL 默认不开启审计，需通过自定义 \textbf{Option Group} 添加 \texttt{MARIADB\_AUDIT\_PLUGIN} 插件
  \item 插件的 \texttt{SERVER\_AUDIT\_EVENTS} 至少要包含能记录 SQL 的事件类型（如 \texttt{QUERY} 或 \texttt{QUERY\_DML}）；若只配 \texttt{CONNECT} 则看不到 SQL
  \item 在 RDS 实例的 \textbf{Modify} 页面同时完成两件事：\\
  （1）关联自定义 Option Group \quad（2）勾选 Audit log 导出至 CloudWatch
  \item 日志组路径固定为 \texttt{/aws/rds/instance/<实例名>/audit}
  \item 审计日志为 CSV 格式，第 3 个字段为 \textbf{用户名}，可用 CloudWatch Insights 解析查询
\end{enumerate}
\end{tcolorbox}

\vspace{2em}
\noindent\rule{\textwidth}{0.4pt}
\begin{center}
\textcolor{gray}{\small 本文档基于 AWS Jam Challenge 实操经验整理 \quad|\quad Amazon RDS + CloudWatch}
\end{center}

\end{document}
